\section{Related Work}
\label{sec:related}

\textbf{Model Merging and Parameter Ensembling.}
The practice of combining multiple trained neural networks into a single model without full re-training has gained significant traction. Model Soups \cite{wortsman22} showed that simply averaging the weights of models fine-tuned with different hyperparameters can improve accuracy and robustness. Git Re-Basin \cite{ainsworth22} addressed the permutation obstacle by aligning the neuron representations of models before averaging. ZipIt \cite{stoica23} extended this alignment to models trained on disparate tasks by identifying and merging redundant features. Task Arithmetic \cite{ilharco22, wang22} and TIES-Merging \cite{yadav23} use task vectors (defined as the difference between fine-tuned and pre-trained weights) to perform linear arithmetic in parameter space, enabling task addition and deletion. While these methods are powerful for weight-space operations, they operate in flat Euclidean spaces and do not preserve the algebraic properties of representation-space projection operators.

\textbf{Dynamic Routing and Activation Blending.}
Rather than statically merging model parameters, dynamic routing methods adjust ensembling coefficients sample-wise during the forward pass. Mixture-of-Experts (MoE) models \cite{jacobs91, jordan94, shazeer17} route tokens to specialized expert sub-networks using learned gating functions, which has scaled models to trillions of parameters \cite{fedus21, jiang24}. Recently, Sample-wise Activation Blending (SABLE) \cite{dauner23} proposed to dynamically blend intermediate activations from different models using sample-specific weights. While SABLE achieves impressive efficiency, it linearly blends activations without imposing manifold constraints, leaving it susceptible to coordinate collapse when task representations reside in non-orthogonal subspaces.

\textbf{Learning-Theoretic Regularization.}
A growing body of work attempts to bring theoretical guarantees to model ensembling. Prior works, such as Rademacher-Bounded Polynomial Merging \cite{maurer04}, restrict ensembling trajectories to low-degree polynomials to control hypothesis class capacity via spectrally-normalized Rademacher complexity bounds. PAC-ZCA \cite{mcallester99, catoni07, donini20} optimizes a temperature-only Gibbs routing policy using parameter-space PAC-Bayesian generalization bounds to provide stable risk guarantees on calibration datasets. However, these methods either use stateless heuristics that suffer from transductive overfitting, or enforce physical trajectory constraints without a unified geometric framework that directly links the optimization objective to the out-of-sample generalization of merged networks under severe noise.

\textbf{Grassmannian Manifolds in Deep Learning.}
The Grassmannian manifold $\mathcal{G}(d, D)$, representing the space of all $d$-dimensional subspaces in $\mathbb{R}^D$, is a fundamental concept in matrix manifold optimization \cite{absil08, edelman98}. Riemannian optimization and geodesic statistics on the Grassmannian have been applied to subspace tracking, computer vision \cite{fletcher04, pennec06, mu20}, and federated learning \cite{shaji21}. In federated learning, some methods communicate local model subspaces rather than raw parameters to reduce bandwidth and protect privacy. Our work is fundamentally distinct: we leverage Grassmannian geodesics and Riemannian-geometric logarithm and exponential mappings to dynamically blend task-specific projection bases sample-wise inside deep neural network routing layers, providing a mathematically rigorous solution to representation-space coordinate collapse.
